\par\indent\SourceLeaf{229}%
Die \textit{zweite} Methode stellt die Koordinaten als Funktionen zweier Variablen \(p,q\) dar. Die Differentiation dieser Funktionen ergebe
\[
\begin{aligned}
\mathrm{d}x&=a\mathrm{d}p+a'\mathrm{d}q\\
\mathrm{d}y&=b\mathrm{d}p+b'\mathrm{d}q\\
\mathrm{d}z&=c\mathrm{d}p+c'\mathrm{d}q
\end{aligned}
\]
\noindent Setzen wir diese Werte in die oben angegebene Formel ein, so erhalten wir
\[
(aX+bY+cZ)\mathrm{d}p+(a'X+b'Y+c'Z)\mathrm{d}q=0
\]
\noindent Da diese Gleichung unabhängig von den Werten der Differentiale \(\mathrm{d}p,\mathrm{d}q\) gelten muss, folgt unmittelbar
\[
aX+bY+cZ=0,\qquad a'X+b'Y+c'Z=0
\]
\noindent Daraus schließen wir, dass \(X,Y,Z\) proportional zu den folgenden Größen sein müssen
\[
bc'-cb',\quad ca'-ac',\quad ab'-ba'
\]
\par\Needspace{10\baselineskip}\indent Setzen wir also der Kürze halber
\[
\BareRoot{((bc'-cb')^2+(ca'-ac')^2+(ab'-ba')^2)}=\Delta
\]
\noindent so gilt entweder
\[
X=\frac{bc'-cb'}{\Delta},\quad Y=\frac{ca'-ac'}{\Delta},\quad Z=\frac{ab'-ba'}{\Delta}
\]
\noindent oder
\[
X=\frac{cb'-bc'}{\Delta},\quad Y=\frac{ac'-ca'}{\Delta},\quad Z=\frac{ba'-ab'}{\Delta}
\]
\par\indent Zu diesen beiden allgemeinen Methoden kommt eine \textit{dritte} hinzu, bei der eine der Koordinaten, etwa \(z\), als Funktion der beiden übrigen \(x,y\) dargestellt wird. Diese Methode ist offenbar ein Spezialfall der ersten oder der zweiten Methode. Setzen wir hier
\[
\mathrm{d}z=t\mathrm{d}x+u\mathrm{d}y
\]
\noindent so gilt entweder
\[
X=\frac{-t}{\BareRoot{(1+tt+uu)}},\quad Y=\frac{-u}{\BareRoot{(1+tt+uu)}},\quad Z=\frac{1}{\BareRoot{(1+tt+uu)}}
\]
\noindent oder
\[
X=\frac{t}{\BareRoot{(1+tt+uu)}},\quad Y=\frac{u}{\BareRoot{(1+tt+uu)}},\quad Z=\frac{-1}{\BareRoot{(1+tt+uu)}}
\]
